% --------------------- NO TOCAR ----------------------
\documentclass[a4paper, 10pt, twocolumn]{report}
\input plantillaTeXpracticasRL.sty % Incluir fichero de estilo del máster
% -----------------------------------------------------



% MODIFICA ESTOS DATOS
% Incluir el nombre de la obra y los nombres completos de los autores del trabajo, junto con el correo electrónico um.es para cada autor.
\title{Titulo del Trabajo: Aprendizaje por Refuerzo\footnote{Comentario que se quiera, si se quiere.}}
\author{
Autor1 	\\
\texttt{autor1@um.es} \\[0.25cm]
Autor2 	\\
\texttt{autor2@um.es} \\[0.25cm]
Autor3 	\\
\texttt{autor3@um.es}}



\begin{document}

% Carátula
\maketitle % Puedes modificar la carátula usando el entorno titlepage.

% Resumen
\input{resumen}

% Índice
\tableofcontents

% Trabajos
\input{bandido}
\input{entornos_complejos}

% Bibliografía
\begin{thebibliography}{9} % El parámetro 9 indica el ancho máximo de la etiqueta, aquí "9" es un ejemplo
	\setlength{\itemsep}{0pt} % No tocar. Reduce aún más el espacio entre ítems 
  	\setlength{\parsep}{0pt}    % No tocar. Elimina el espacio adicional entre párrafos dentro de un ítem
  	
    \bibitem{ref1} % Para citar usar \cite{ref1}
    Autor, A. (Año). \textit{Título del artículo}. Nombre de la revista, Volumen(Número), páginas.

    \bibitem{ref2} % Para citar usar \cite{ref2}
    Autor, B. (Año). \textit{Título del libro}. Editorial.
    
\end{thebibliography}


Listar todas las referencias utilizadas en el trabajo, siguiendo un formato bibliográfico estándar de \LaTeX\/ o de alguna asociación (APA, IEEE, etc.). La información obligatoria para que tipo de publicación la puede consultar en \url{https://bibtex.eu/fields/} Se recomienda utilizar un gestor de referencias como Bib\TeX\/ para facilitar la organización de las citas solo si va a poner mucha bibliografía. 


\end{document}